\documentclass[a4paper,10pt]{article}
\usepackage[utf8]{inputenc}
\usepackage[T1]{fontenc}
\usepackage[french]{babel}
\usepackage[margin=1.5cm]{geometry}
\usepackage{xcolor}
\usepackage{tcolorbox}
\usepackage{enumitem}
\usepackage{multicol}
\usepackage{lmodern}
\usepackage{tabularx}
\usepackage{booktabs}
\usepackage{listings}

\renewcommand{\familydefault}{\sfdefault}
\pagestyle{empty}
\setlist{nosep, leftmargin=1.5em}

\definecolor{navy}{RGB}{0,51,102}
\definecolor{teal}{RGB}{0,120,120}
\definecolor{crimson}{RGB}{180,0,0}
\definecolor{green}{RGB}{0,120,50}
\definecolor{lightbg}{RGB}{245,248,252}
\definecolor{orange}{RGB}{200,100,0}
\definecolor{purple}{RGB}{100,0,160}
\definecolor{codebg}{RGB}{28,28,36}
\definecolor{codecyan}{RGB}{100,200,220}
\definecolor{codeorange}{RGB}{255,165,0}
\definecolor{codegreen}{RGB}{100,200,100}
\definecolor{codemagenta}{RGB}{220,100,220}

\lstdefinestyle{html}{
  backgroundcolor=\color{codebg},
  basicstyle=\ttfamily\small\color{white},
  keywordstyle=\color{codecyan}\bfseries,
  stringstyle=\color{codeorange},
  commentstyle=\color{gray}\itshape,
  frame=none, breaklines=true,
  language=HTML,
  morekeywords={header, nav, main, section, article, aside, footer, figure, figcaption},
}

\lstdefinestyle{css}{
  backgroundcolor=\color{codebg},
  basicstyle=\ttfamily\small\color{white},
  keywordstyle=\color{codemagenta}\bfseries,
  stringstyle=\color{codeorange},
  commentstyle=\color{gray}\itshape,
  frame=none, breaklines=true,
  extendedchars=true,
  inputencoding=utf8,
}

\lstdefinestyle{js}{
  backgroundcolor=\color{codebg},
  basicstyle=\ttfamily\small\color{white},
  keywordstyle=\color{codecyan}\bfseries,
  stringstyle=\color{codeorange},
  commentstyle=\color{gray}\itshape,
  frame=none, breaklines=true,
  language=Java,
}

\lstdefinestyle{php}{
  backgroundcolor=\color{codebg},
  basicstyle=\ttfamily\small\color{white},
  keywordstyle=\color{codegreen}\bfseries,
  stringstyle=\color{codeorange},
  commentstyle=\color{gray}\itshape,
  frame=none, breaklines=true,
  language=PHP,
}

\tcbset{basebox/.style={
  colback=lightbg, colframe=navy, coltitle=white,
  fonttitle=\bfseries, arc=2mm, boxrule=0.5mm,
  left=3mm, right=3mm, top=2mm, bottom=2mm, after skip=6pt
}}

\newcommand{\key}[1]{\textbf{\textcolor{navy}{#1}}}
\newcommand{\warn}[1]{\textbf{\textcolor{crimson}{#1}}}

\begin{document}

\begin{center}
  \colorbox{orange}{\parbox{\dimexpr\textwidth-2\fboxsep}{
    \centering\vspace{3pt}
    {\LARGE\bfseries\textcolor{white}{CRÉATION WEB — HTML, CSS, JavaScript, PHP}}\\
    {\normalsize\textcolor{white}{BTS SIO — Maël Mignard}}
    \vspace{3pt}
  }}
\end{center}
\vspace{6pt}

\begin{multicols}{2}

% ============================================================
\begin{tcolorbox}[basebox, colframe=orange, title=1. HTML — Structure d'une page]
\begin{lstlisting}[style=html]
<!DOCTYPE html>
<html lang="fr">
<head>
  <meta charset="UTF-8">
  <meta name="viewport"
        content="width=device-width">
  <title>Ma page</title>
  <link rel="stylesheet" href="style.css">
</head>
<body>
  <header>
    <nav>
      <a href="index.html">Accueil</a>
    </nav>
  </header>
  <main>
    <h1>Titre principal</h1>
    <p>Paragraphe de texte.</p>
    <img src="img.jpg" alt="Description">
  </main>
  <footer>
    <p>Pied de page</p>
  </footer>
  <script src="script.js"></script>
</body>
</html>
\end{lstlisting}
\end{tcolorbox}

% ============================================================
\begin{tcolorbox}[basebox, colframe=orange, title=2. HTML — Balises importantes]
\begin{tabular}{ll}
\toprule
\textbf{Balise} & \textbf{Rôle} \\
\midrule
\texttt{<h1>...<h6>} & Titres (hiérarchie) \\
\texttt{<p>} & Paragraphe \\
\texttt{<a href="">} & Lien hypertexte \\
\texttt{<img src="" alt="">} & Image \\
\texttt{<ul> / <ol>} & Liste non/ordonnée \\
\texttt{<li>} & Élément de liste \\
\texttt{<div>} & Conteneur bloc (générique) \\
\texttt{<span>} & Conteneur inline \\
\texttt{<form>} & Formulaire \\
\texttt{<input type="">} & Champ de saisie \\
\texttt{<button>} & Bouton \\
\texttt{<table>} & Tableau HTML \\
\texttt{<strong>} & Texte important (gras) \\
\texttt{<em>} & Texte en emphase (italic) \\
\bottomrule
\end{tabular}
\end{tcolorbox}

% ============================================================
\begin{tcolorbox}[basebox, colframe=teal, title=3. CSS — Sélecteurs et propriétés]
\begin{lstlisting}[style=css]
/* Selecteur de balise */
h1 { color: navy; font-size: 2em; }

/* Selecteur de classe */
.card { background: white; padding: 1rem; }

/* Selecteur d'id */
#header { width: 100%; }

/* Combinaisons */
nav > a { color: blue; }
div:hover { opacity: 0.8; }
input::placeholder { color: gray; }

/* Boite */
.box {
  margin: 10px;       /* exterieur */
  border: 1px solid;  /* bordure */
  padding: 15px;      /* interieur */
  width: 300px;
}
\end{lstlisting}
\end{tcolorbox}

% ============================================================
\begin{tcolorbox}[basebox, colframe=teal, title=4. CSS — Flexbox \& Grid]
\begin{lstlisting}[style=css]
/* FLEXBOX */
.container {
  display: flex;
  flex-direction: row;     /* ou column */
  justify-content: center; /* axe principal */
  align-items: center;     /* axe secondaire */
  gap: 1rem;
  flex-wrap: wrap;
}

/* GRID */
.grille {
  display: grid;
  grid-template-columns: repeat(3, 1fr);
  grid-gap: 20px;
}

/* RESPONSIVE */
@media (max-width: 768px) {
  .container {
    flex-direction: column;
  }
}
\end{lstlisting}
\key{Flexbox} : 1D (ligne ou colonne) \\
\key{Grid} : 2D (lignes ET colonnes)
\end{tcolorbox}

% ============================================================
\begin{tcolorbox}[basebox, colframe=navy, title=5. JavaScript — Manipulation du DOM]
\begin{lstlisting}[style=js]
// Selectionner des elements
const btn = document.getElementById("btn");
const items = document.querySelectorAll(".item");

// Modifier le contenu
btn.textContent = "Cliquez";
btn.innerHTML = "<strong>OK</strong>";

// Modifier le style
btn.style.color = "red";
btn.classList.add("active");
btn.classList.remove("hidden");

// Evenements
btn.addEventListener("click", function() {
  alert("Bouton clique !");
});

// Creer et ajouter un element
const p = document.createElement("p");
p.textContent = "Nouveau paragraphe";
document.body.appendChild(p);
\end{lstlisting}
\end{tcolorbox}

% ============================================================
\begin{tcolorbox}[basebox, colframe=navy, title=6. JavaScript — Fetch API (AJAX)]
\begin{lstlisting}[style=js]
// GET : recuperer des donnees
fetch("https://api.example.com/data")
  .then(response => response.json())
  .then(data => {
    console.log(data);
  })
  .catch(error => console.error(error));

// POST : envoyer des donnees
fetch("/api/users", {
  method: "POST",
  headers: { "Content-Type": "application/json" },
  body: JSON.stringify({ nom: "Mael", age: 20 })
})
  .then(r => r.json())
  .then(data => console.log(data));

// Avec async/await (plus lisible)
async function getData() {
  const r = await fetch("/api/data");
  const data = await r.json();
  return data;
}
\end{lstlisting}
\end{tcolorbox}

% ============================================================
\begin{tcolorbox}[basebox, colframe=green, title=7. PHP — Bases côté serveur]
\begin{lstlisting}[style=php]
<?php
// Variables
$nom = "Mael";
$age = 20;
$majeur = $age >= 18;

// Affichage
echo "Bonjour $nom";
echo "<br>";

// Condition
if ($majeur) {
    echo "Majeur";
} else {
    echo "Mineur";
}

// Tableaux
$notes = [12, 15, 18];
foreach ($notes as $note) {
    echo $note . "<br>";
}

// Fonctions
function somme($a, $b) {
    return $a + $b;
}
echo somme(3, 4); // 7
?>
\end{lstlisting}
\end{tcolorbox}

% ============================================================
\begin{tcolorbox}[basebox, colframe=green, title=8. PHP — Formulaires et BDD]
\begin{lstlisting}[style=php]
<?php
// Récupérer données formulaire (SÉCURISÉ)
$nom = htmlspecialchars($_POST["nom"]);
$email = filter_var($_POST["email"],
         FILTER_VALIDATE_EMAIL);

// Connexion MySQL (PDO - recommandé)
$pdo = new PDO(
    "mysql:host=localhost;dbname=mabase",
    "user", "password"
);

// Requête préparée (anti-SQLi)
$stmt = $pdo->prepare(
    "SELECT * FROM users WHERE email = ?"
);
$stmt->execute([$email]);
$user = $stmt->fetch(PDO::FETCH_ASSOC);

// Sessions
session_start();
$_SESSION["user"] = $nom;

// Cookies
setcookie("token", "abc123", time()+3600);
?>
\end{lstlisting}
\warn{Toujours utiliser des requêtes préparées} (PDO/MySQLi) pour éviter les injections SQL !
\end{tcolorbox}

% ============================================================
\begin{tcolorbox}[basebox, colframe=crimson, title=9. Sécurité Web — À ne jamais oublier]
\begin{itemize}
  \item \warn{XSS} : toujours utiliser \texttt{htmlspecialchars()} sur les données affichées
  \item \warn{SQLi} : toujours utiliser des requêtes préparées (PDO)
  \item \warn{CSRF} : utiliser des tokens de formulaire
  \item \warn{Uploads} : valider le type MIME, ne pas exécuter les fichiers uploadés
  \item \warn{Mots de passe} : hasher avec \texttt{password\_hash()} / \texttt{password\_verify()}
  \item HTTPS partout (certificat TLS)
  \item Valider les données côté serveur (jamais juste côté client)
\end{itemize}
\end{tcolorbox}

% ============================================================
\begin{tcolorbox}[basebox, colframe=purple, title=10. WordPress — Points clés]
\key{CMS} : Content Management System \\
\key{Thème} : apparence du site \\
\key{Plugin} : fonctionnalités supplémentaires \\
\key{Pages} vs \key{Articles} : statique vs dynamique \\[4pt]
\textbf{Structure fichiers WordPress :}
\begin{itemize}
  \item \texttt{wp-content/themes/} : thèmes
  \item \texttt{wp-content/plugins/} : plugins
  \item \texttt{wp-config.php} : config BDD
\end{itemize}
\vspace{4pt}
\key{WP-CLI} : gérer WordPress en ligne de commande \\
\key{BDD} : MySQL, tables préfixées \texttt{wp\_}
\end{tcolorbox}

\end{multicols}

\vspace{4pt}
\begin{tcolorbox}[basebox, colframe=orange, title=Points clés à retenir pour l'examen]
\begin{multicols}{2}
\begin{itemize}
  \item HTML = structure, CSS = style, JS = comportement, PHP = serveur
  \item \texttt{<!DOCTYPE html>} : déclaration obligatoire en HTML5
  \item Balises sémantiques : \texttt{<header>}, \texttt{<nav>}, \texttt{<main>}, \texttt{<footer>}
  \item \texttt{@media (max-width: 768px)} : breakpoint responsive mobile
  \item Flexbox : \texttt{display:flex} — 1D ; Grid : \texttt{display:grid} — 2D
  \item \texttt{htmlspecialchars()} = protection contre XSS en PHP
  \item Requêtes préparées PDO = protection contre SQLi
  \item \texttt{\$\_POST} = données formulaire (POST), \texttt{\$\_GET} = URL (GET)
  \item \texttt{session\_start()} avant tout dans un fichier PHP avec sessions
  \item JS : \texttt{querySelector}, \texttt{addEventListener}, \texttt{fetch} sont les 3 essentiels
\end{itemize}
\end{multicols}
\end{tcolorbox}

\end{document}
